\section{Life of the Spirit}

\textit{Originally published in November, 2005 for a private seminar studying non-dual writings.}

\begin{quotex}
The race of the spirit predominant in a person, in a people, in a community, is given by the orienting characteristic assumed toward the sacred and the divine, life and death, destiny, the world.

\textit{Noi possediamo una verità, non una retorica.}

We possess a truth, not a rhetoric.
\flright{\textsc{Julius Evola}}
\end{quotex}


The Traditional view is that man is a tripartite being: body, soul, and spirit. The body represents our phenomenal existence, the spirit our noumenal existence, and the soul --- though phenomenal --- is suspended between them.

The soul is the animating, or integrating, principle of the body. In a healthy personality, the soul dominates the body, and the soul is the expression of the spirit, reflecting its fundamental and orienting principles.

In a disordered personality, the center of gravity of consciousness may be situated in the body, and the existence of higher states is only dimly present to consciousness, or else totally denied. In such a state, then only the desires, pleasures, and pains of the body are of any significance.

If consciousness is centered in the soul, then life is dominated by either emotional states, or by a rationality that denies any sort of higher knowledge, wisdom, or gnosis.

If consciousness is centered in spirit, then the center is the noumenal self, the silent witness, the ``unmoved mover'', that is unaffected by anything in the phenomenal realm, yet is the source of all that manifests.

It is important, then, to realize that anything in manifestation is not by accident, but has its source, rationale, and \emph{raison d'être} in spirit. In particular, one's gender, race, ethnicity, character, qualities, even historical circumstances are a reflection of this noumenal Self.

This is the esoteric understanding of ``Karma'': karma is not the result of some accident, or event, or act, nor is it imposed from the outside, so to speak, perhaps as some sort of punishment or reward. Rather, it arises out of the very nature of the Self and is its own self-expression in phenomena.

As noumenon, the Self is beyond Being, and therefore partakes of Non-being, that is, its unrealized possibilities. As life is the actualization of possibilities, this brings us to the esoteric meaning of ``Dharma''. Dharma is our destiny: our task, then, is to manifest the Self in phenomena, or to say it more plainly, to ``become who we are''.

To the exoteric mind, this is an outrage. It regards these differentiations as unfair and arbitrary. To be born a man or woman, a prince or pauper, good or evil is simply a genetic accident. It misunderstands the progressive unfolding of the higher Self in history as the evolution of the lower self to the higher self.

In this diagram\footnote{The image is no longer available in the original site. It is likely it is referring to the stereographical projection.}, we can regard point N as the noumenal self and the plane as the phenomenal world. Every point on the sphere is projected onto a point on the plane, with its origin at point N; for example, Q maps to P', and R to P. S, the lower self, is the tangent point where the sphere meets the plane. Note, however, that point N itself does not map anywhere onto the plane, even though it is the source of all the projections.\footnote{The circle on the plane represents the lifespan of the Self in a particular plane of existence --- more on this in due course.}

Fate, destiny, freedom: these three concepts can now be understood. Fate is necessity on the phenomenal plane: the present is determined by the past. Destiny is necessity in the noumenal sense; it is our dharma, what we are called to be(come). As such it differs from fate, which is necessity compelled from the outside, whether in the corporal sense as impulses from other bodies, or even in the soul sense, as impulses from psychological and other mental complexes.

Destiny in its fullest expression is self-determined, the Self-expression emanating from our deepest source. This is experienced as the mysterious source of creativity and will. Note that this is the esoteric meaning of ``Freedom'' --- it does not mean that we could have chosen otherwise; rather it means that our choices emanate from our deepest self and are therefore self-determined --- not ``other'' determined. If our choices arise from a conflict of other forces, whether spiritual and material, then such a choice cannot be said to be fully free.

Evola, in this book, intends to reveal the spiritual forces that drive world-historical events. These forces are analyzed both as micro forces acting on the individual and as macro forces acting on communities and nations. Referring back to the diagram above, most works of social analysis confine themselves to the plane of phenomena, and analyze the material, economic, and psychological forces to account for world events, as though such an analysis exhausts all possibilities. For us, instead, we will focus on the ``third'' dimension --- the spiritual dimension --- and look at the occult, or hidden, forces that impinge upon the plane of phenomena.


\flrightit{Posted on 2011-04-12 by Cologero}

\begin{center}* * *\end{center}

\begin{footnotesize}
\begin{sffamily}
\texttt{HOO on 2011-04-13 at 23:10 said:}

One of the best posts so far. O, sweet metaphysics.

\end{sffamily}
\end{footnotesize}

